\documentclass[12pt, a4paper]{article}
\usepackage[icelandic]{babel}
\usepackage[utf8]{inputenc}
\usepackage[T1]{fontenc}
\usepackage{amsmath, amssymb}
\usepackage{geometry}
\usepackage{booktabs}
\usepackage{hyperref}
\geometry{margin=2.5cm}

\title{Aðferðir -- Vatnafræðigreining Ölfusár við Selfoss}
\author{}
\date{}

\begin{document}
\maketitle
\tableofcontents
\newpage

% ═══════════════════════════════════════════════════════════════════════════════
\section{Gögn og tímabil}
% Skrift: 00\_hreinsa\_gogn.py
% ═══════════════════════════════════════════════════════════════════════════════

Gögnin eru sótt úr LamaH-Ice gagnasafninu fyrir Ölfusá við Selfoss (stöð ID~98).
Notaðar eru dagleg mælingar á rennsli [m$^3$/s], úrkomu [mm/dag] og hitastigi [$^\circ$C]
úr CARRA endurgreiningu, yfir vatnafræðilegt tímabil 1993--2023 (30~ár,
1.~október~1993 -- 30.~september~2023).

Gögn um rennsli og veðurfar eru síuð yfir tímabilið og úthlutað í tvær hreinar
CSV-skrár sem allar niðurstreymis skriptur nota. Neikvæð rennslisgildi eru sett
sem vantar (\texttt{NaN}). Úrkoma er sótt úr CARRA endurgreiningu og er ERA5-Land
gögn notuð sem varabúnaður þar sem CARRA gögn vantar.

% ═══════════════════════════════════════════════════════════════════════════════
\section{Lýsing vatnasviðs}
% Skrift: 01\_lysing\_vatnasviðs.py
% ═══════════════════════════════════════════════════════════════════════════════

Eiginleikar vatnasviðsins eru sóttir úr \texttt{Catchment\_attributes.csv}
gagnaskránni og kynntir í töfluformi og súluriti. Meðal þeirra eiginleika sem
eru kynntir eru flatarmál, meðalhæð, miðgildi hæðar, hæðarbil, meðalhalli,
löngungarstuðull, lækjaþéttleiki, landnýtingarhlutföll (jökull, bert land,
kjarr, votlendi, landbúnaður, skógur, stöðuvatn, þéttbýli), veðurfar og
grunnvatnshlutfall.

% ═══════════════════════════════════════════════════════════════════════════════
\section{Árstíðasveifla}
% Skrift: 02\_arstidarsveifla.py
% ═══════════════════════════════════════════════════════════════════════════════

Mánaðarleg meðaltöl (climatology) fyrir rennsli, úrkomu og hitastig eru reiknuð
með meðaltali yfir öll 30~ár og sett fram á þríhluta mynd. Staðalfrávik er
reiknað til að sýna náttúrulega sveiflukenningu.

Mánaðarsummur úrkomu eru reiknaðar á hvern mánuð á hverju ári
(með \texttt{resample("ME").sum()}) og síðan tekið meðaltal og staðalfrávik
þessara 30 mánaðarsumma til að gefa rétta lýsingu á mánaðarlegri úrkomudreifingu.

% ═══════════════════════════════════════════════════════════════════════════════
\section{Mat á grunnvatnsframlagi}
% Skrift: 03\_grunnvatnsframlag.py
% ═══════════════════════════════════════════════════════════════════════════════

Þrjár aðferðir eru notaðar til að meta grunnvatnsframlag í heildarrennsli.

\subsection{Lyne-Hollick stafræn sía}

Síunni er beitt til að skilja á milli grunnvatnsrennslis (baseflow) og
skyndirennslis (quickflow) \cite{Ladson2013}. Jafnan er:

\begin{equation}
    f_t = \alpha \cdot f_{t-1} + \frac{1+\alpha}{2}\bigl(Q_t - Q_{t-1}\bigr)
    \label{eq:lyne-hollick}
\end{equation}

þar sem $f_t$ er skyndirennsli á degi~$t$, $Q_t$ er heildarrennsli og
$\alpha = 0{,}925$ er síufæribreyta \cite{Ladson2013}.
Síunni er beitt þrjú skipti (fram--aftur--fram) til að fjarlægja fasaflökt.
Grunnvatn er metið sem $b_t = Q_t - f_t$, þar sem $b_t$ er klippt í bilið
$[0,\, Q_t]$.

\subsection{Grunnvatnshlutfall (Baseflow Index)}

Grunnvatnshlutfall (BFI) er reiknað sem:

\begin{equation}
    \mathrm{BFI} = \frac{\displaystyle\sum_{t} b_t}{\displaystyle\sum_{t} Q_t}
    \label{eq:bfi}
\end{equation}

þar sem gildi nálægt~1 þýðir að grunnvatn ræður ríkjum, en gildi nálægt~0
þýðir að skyndirennsli er ríkjandi. BFI er einnig reiknað fyrir hvert
vatnafræðilegt ár til að meta breytileika milli ára.

\subsection{Recession-greining}

Samfelldir samdráttarhlutАР eru greindir þar sem rennsli lækkar samfellt og
enginn úrkoma er til staðar ($P < 0{,}5$~mm/dag í að minnsta kosti 5~daga).
Recession constant~$k$ er metin með línulegri aðhvarfsgreiningu á $\ln(Q)$
sem fall af tíma:

\begin{equation}
    Q(t) = Q_0 \cdot e^{-kt}
    \label{eq:recession}
\end{equation}

þar sem $Q_0$ er rennsli við upphaf samdráttar og $k$ er recession constant
[dag$^{-1}$]. Miðgildi $k$ yfir öll greind brot er notað sem dæmigerð gildi.
Úr $k$ eru einnig leiddar geymslufastin $\tau = 1/k$ og halftíminn
$t_{1/2} = \ln 2 / k$.

% ═══════════════════════════════════════════════════════════════════════════════
\section{Vatnsjafnaðargreining}
% Skrift: 04\_grunnliking.py
% ═══════════════════════════════════════════════════════════════════════════════

Grunnlíking vatnafræðinnar er notuð til að meta vatnsbúskap vatnasviðsins:

\begin{equation}
    P = Q + ET + \Delta S
    \label{eq:water-balance}
\end{equation}

þar sem $P$ er úrkoma, $Q$ er afrennsli, $ET$ er uppgufun (úr CARRA
endurgreiningu) og $\Delta S$ er breyting á geymslu, reiknuð sem leifð.
Rennsli er umreiknað úr m$^3$/s í mm/dag með:

\begin{equation}
    Q_{\mathrm{mm}} = Q_{\mathrm{m^3/s}} \cdot \frac{86\,400}{A \times 10^6}
    \times 1000
    \label{eq:q-conversion}
\end{equation}

þar sem $A$ er flatarmál vatnasviðsins í km$^2$. Jákvætt $\Delta S$ þýðir
aukningu í geymslu (snjór, jökull, grunnvatn) en neikvætt $\Delta S$ þýðir
geymslutæmingu, t.d.\ vegna jökulbræðslu. Greiningin er gerð bæði á daglegum
gögnum og á árslegum vatnafræðilegum árum (október--september).

% ═══════════════════════════════════════════════════════════════════════════════
\section{Langæislína (Flow Duration Curve)}
% Skrift: 05\_langaeslina.py
% ═══════════════════════════════════════════════════════════════════════════════

Öll daglegt rennsli er raðað í lækkandi röð og yfirfallstíðni reiknuð með
Weibull plotting position:

\begin{equation}
    P(i) = \frac{i}{N+1} \times 100\%
    \label{eq:weibull}
\end{equation}

þar sem $i$ er röðunarnúmer (1~=~hæsta gildi) og $N$ er heildarfjöldi gilda.
Lykilgildi Q5, Q50 og Q95 eru fundin með línulegri brúun, og hlutfallið
$Q5/Q95$ er notað sem mælikvarði á sveiflukenni rennslis. Langæislínan er
einnig teiknuð á árstíðabundnum grunni (vor, sumar, haust, vetur) til að
sýna hvernig dreifingin breytist eftir árstímum.

% ═══════════════════════════════════════════════════════════════════════════════
\section{Flóðagreining}
% Skrift: 06\_flodagreining.py
% ═══════════════════════════════════════════════════════════════════════════════

\subsection{Flóðatíðni eftir mánuðum}

Árstoppar eru greindir eftir vatnafræðilegu ári (október--september) og
flokkaðir eftir mánuðum til að meta árstíðabundin mynstur í flóðatíðni.

\subsection{Líkindadreifingargreining}

Þrjár líkindadreifingar eru passar á árstoppa með Gringorten plotting
positions \cite{Gringorten1963}:

\begin{equation}
    F(i) = \frac{i - 0{,}44}{N + 0{,}12}
    \label{eq:gringorten}
\end{equation}

þar sem $F(i)$ er uppsöfnuð líkindafall og endurkomutími er $T = 1/(1-F)$.
Dreifinganna þrjár sem eru prófaðar eru:

\begin{itemize}
    \item \textbf{Gumbel (EV Type~I):}
    \begin{equation}
        F(x) = \exp\!\bigl(-e^{-(x-\mu)/\beta}\bigr), \quad
        \beta = \frac{\sigma\sqrt{6}}{\pi}, \quad
        \mu = \bar{x} - 0{,}5772\,\beta
        \label{eq:gumbel}
    \end{equation}

    \item \textbf{Log Normal 3 stika (LN3):}
    $\log(X - \xi) \sim \mathcal{N}(\mu_{\ln},\, \sigma_{\ln}^2)$
    þar sem neðri mörk $\xi$ eru fundin með lágmörkun RMSE á CDF-passan
    og stikamöt framkvæmd með MLE.

    \item \textbf{Log Pearson 3 stika (LP3):}
    $y = \log_{10}(Q)$ fylgir Pearson Type~III dreifingu með stika
    $\bar{y}$, $s$ og skekju~$g$ (sample skewness).
\end{itemize}

Besta dreifing er valin með lægstu samlægri RMSE á heildarpassun og
halapassun (efri 25\% toppa). Endurkomutímar Q10, Q50 og Q100 eru reiknuð
með 90\% öryggisbilum, metin með bootstrap-aðferð
($N_{\mathrm{boot}} = 2000$ endurteiknanir, 5.\ og 95.\ hundraðshlutfall).

% ═══════════════════════════════════════════════════════════════════════════════
\section{Leitnigreining}
% Skrift: 07\_leitnigreining.py
% ═══════════════════════════════════════════════════════════════════════════════

Tvær tölfræðilegar aðferðir eru notaðar samhliða til að greina leitni í
rennsli, úrkomu og hitastigi.

\subsection{Modified Mann-Kendall próf}

Modified Mann-Kendall próf \cite{Hamed1998} leiðréttir fyrir sjálffylgni í
tímaröðum, sem er algeng í vatnafræðilegum gögnum. Marktækni er metin við
$p < 0{,}05$. Prófið er beitt á árlegar, mánaðarlegar og árstíðabundnar
tímaraðir.

\subsection{Theil-Sen hallatala}

Theil-Sen hallatalan \cite{Sen1968, Theil1950} er öflugt mat á leitni sem er
ónæmt fyrir útlögum. Hallatalan er metin sem miðgildi allra parpara
hlutfallslegra breytinga:

\begin{equation}
    \hat{\beta} = \operatorname{median}\!\left(\frac{y_j - y_i}{x_j - x_i}
    \right), \quad i < j
    \label{eq:theil-sen}
\end{equation}

Heildarbreyting yfir tímabilið er reiknuð sem hlutfallsleg breyting á milli
upphafs og loka aðhvarfslínunnar. Leitni er greind á þremur stigum: árlegum,
mánaðarlegum (fyrir hvern mánuð) og árstíðabundnum (vor, sumar, haust, vetur).

% ═══════════════════════════════════════════════════════════════════════════════
\section{Greining á rennslisatburði}
% Skrift: 08\_rennslisatburdur.py
% ═══════════════════════════════════════════════════════════════════════════════

Stærsta flóð á mælingatímabilinu (21.~desember~2006,
$Q_{\max} = 1931$~m$^3$/s) er greint í smáatriðum sem dæmi um
hlýviðurflóð (rain-on-snow). Í aðdraganda flóðsins var kalt ($-7$ til
$-9\,^\circ$C) og mikill snjór lagðist á vatnasviðið (6.--17.~des.).
Þegar hlýtt rigndi 18.--21.~des.\ ($T > 0\,^\circ$C) bræddu bæði snjór
og rigning saman og olli skyndilegri flóðbylgju.

Lykilstærðir sem metnar eru:

\begin{itemize}
    \item \textbf{Time-to-peak:} tími frá upphafi rennslisaukningar að
    $Q_{\max}$.

    \item \textbf{Recession time:} metið með $k$ úr jöfnu~\eqref{eq:recession}:
    \begin{equation}
        t_{\mathrm{rec}} = \frac{\ln(Q_{\max}/Q_{\mathrm{base}})}{k}
        \label{eq:recession-time}
    \end{equation}
    þar sem $Q_{\mathrm{base}}$ er grunnrennsli fyrir flóðið og $k$ er
    recession constant metin í skrefi~4.

    \item \textbf{Excess rain release:} tími frá lokum meginúrkomu þangað
    til rennsli snýr aftur að $Q_{\mathrm{base}}$, metið með
    $t_{\mathrm{rec}}$.
\end{itemize}

Atburðurinn er settur fram á þríhluta mynd með rennsli, úrkomu og hitastigi
samhliða til að sýna tengsl milli veðurfarlegs og vatnafræðilegs þróunar
flóðsins.

% ═══════════════════════════════════════════════════════════════════════════════
\begin{thebibliography}{9}

\bibitem{Ladson2013}
Ladson, A.~R., Brown, R., Neal, B., og Nathan, R. (2013).
A standard approach to baseflow separation using the Lyne and Hollick filter.
\textit{Australasian Journal of Water Resources}, 17(1), 25--34.

\bibitem{Hamed1998}
Hamed, K.~H. og Ramachandra Rao, A. (1998).
A modified Mann-Kendall trend test for autocorrelated data.
\textit{Journal of Hydrology}, 204(1--4), 182--196.

\bibitem{Sen1968}
Sen, P.~K. (1968).
Estimates of the regression coefficient based on Kendall's tau.
\textit{Journal of the American Statistical Association}, 63(324), 1379--1389.

\bibitem{Theil1950}
Theil, H. (1950).
A rank-invariant method of linear and polynomial regression analysis.
\textit{Proceedings of the Royal Netherlands Academy of Sciences}, 53, 386--392.

\bibitem{Gringorten1963}
Gringorten, I.~I. (1963).
A plotting rule for extreme probability paper.
\textit{Journal of Geophysical Research}, 68(3), 813--814.

\end{thebibliography}

\end{document}
